\chapter{Toward Portable Confidential VM Self-protection with Operating System Supports (Work-in-progress)}
\renewcommand{\thename}{\oblivm}
% \pagenumbering{arabic}
\label{chap:oblivm}

Confidential Computing technologies such as AMD SEV and Intel TDX enable entire virtual machines to be executed within a confidentiality- and integrity-protected environment. However, a malicious hypervisor can still perform \emph{side-channel} attacks on the protected VM to extract information or hijack the control flow.

Previously, side-channel protection mechanisms were discussed in the context of Intel SGX enclaves. Previous works leverage compiler instrumentations to transform program memory access patterns or insert inline checks to detect attack attempts.
Unfortunately, such approaches lack (1) portability due to maintaining complex compiler toolchains and (2) flexibility in implementing diverse defense mechanisms.

On the other hand, with the adoption of unikernels, many privileged primitives, such as the page table and the fault handler, become efficient for the applications to access.
We observe that side-channel defenses could be implemented on top of these primitives, alleviating the need for a separate compiler toolchain.

This paper is a design exploration of operating system-supported confidential VM side-channel defense.
First, we designed a paging scheme that implements ORAM to obfuscate the \emph{physical memory access pattern} of a VM. Next, we introduced a page fault-based instruction emulation primitive. Leveraging this primitive, we created a salting-on-access scheme and an in-register sensitive data protection scheme to defend against the recently discovered ciphertext side-channel and the cache-based data-only attack.

We are implementing our prototype on top of the Unikraft unikernel. The plan is to study the performance impact of such schemes on real-world applications and also to explore different optimizations to improve the performance.


\section{Background}
\subsection{Secure Encrypted Virtualization}

AMD \gls{sev} is a set of architectural extensions to AMD CPUs that enable the deployment of \glspl{cvm}.
\gls{sev} piggy-backs on the AMD \gls{sme}'s support for transparent memory encryption automatically performed by the \gls{sp}.
When \gls{sev} is enabled, each \gls{cvm} is provisioned a unique \gls{vek} that the \gls{sp} uses to encrypt its memory pages. A \gls{cvm}'s \gls{vek} is securely maintained by the CPU, and is indexed by the \gls{cvm}'s unique \gls{asid}.




\hdr{Address translation SEV}
Similar to traditional \glspl{vm}, \gls{cvm} in \gls{sev} also follows \gls{slat} provided by the hardware.
The MMU translates a virtual address accessed by software inside a \gls{cvm} in two levels. The first level is the translation between \gls{gva} and \gls{hpa}, which is . When a \gls{gpf}.

The page table also contains information about whether a page should be encrypted or not, in a bit in the PTE called the \texttt{c}-bit in both the \gls{gpt}'s and the \gls{npt}'s entries.
% To decide whether a page should be encrypted or not, \gls{sev} relies on the in .


% In the \gls{gpt}

\hdr{Security extensions}
Since \gls{sev}'s release, several vulnerabilities have been discovered by the research community.
For instance, the initial release of \gls{sev} does not protect the \gls{cvm} states that are stored in the \gls{vmcb}, allowing a malicious hypervisor to inject register values alter the \gls{cvm} execution or extract sensitive values that are stored in the registers~\cite{morbitzer_extracting_2019}.
Moreover, the lack of protection in the \gls{gpa} to \gls{hpa} mappings enables the hypervisor to replace the mapping or map a single \gls{gpa} to multiple \glspl{hpa}.
In response, AMD introduced two new security extensions.

\Gls{seves} introduce two key security features. First, all the registers stored in the virtual machine control block (VMCS) are automatically encrypted upon exiting a \gls{vm}, to prevent leakages. Second, Non-automatic Asynchronous Exit (NAE) and the \gls{ghcb} are introduced to allow the guest to selectively share states with the hypervisor when required.

The next security extension to \gls{sev} is \gls{sevsnp}, which adds integrity protection to encrypted pages.
It introduces a \gls{rmp} that tracks the metadata (e.g., owner of the page and the mapped \gls{gpa}) for each \gls{hpa}.
The processor then looks up the \gls{rmp} and checks for permission on every memory access.


\subsection{Side channels on CVMs}

\subsubsection{Nested Page Fault Side-channels}
% \begin{table}[t]
    \small
    \centering
    \begin{tabularx}{\columnwidth}{@{}lX@{}}
    \toprule
         \thead{Attack}& \thead{Leaked Pattern}  \\
         \midrule
         \cite{hetzelt2017security,werner2019severest}& System calls \\
         Se\cite{werner2019severest} & Application structure\\
         \cite{morbitzer2018severed,li2019exploiting,wilke2020sevurity,radev2020exploiting,li2021cipherleaks,li2022systematic} & Data accesses during critical window\\
         \cite{li2022systematic} & Function called during critical window \\
         \cite{li2022systematic}& Context switches\\
         \bottomrule
    \end{tabularx}
    \caption{Previously proposed attacks on SEV that requires the NPF controlled-channel}
    \label{tab:npf-attack}
\end{table}
% \begin{table}[t]
%     \small
%     \centering
%     \begin{tabularx}{\columnwidth}{@{}lccX@{}}
%     \toprule
%          \thead{Attack}& \multicolumn{2}{c}{\thead{Leaked Access Pattern}} &\thead{Implication} \\
%          & \thead{Code}  & \thead{Data}& \\ 
%          \midrule
%          \cite{hetzelt2017security,werner2019severest}&\\
%          \cite{werner2019severest} &\\
%          \cite{morbitzer2018severed,li2019exploiting,wilke2020sevurity,radev2020exploiting,li2021cipherleaks,li2022systematic} & \\ 
%          \cite{li2022systematic} & \\
%          \cite{li2022systematic}& \\
%          SEVerest~\cite{werner2019severest}& & &\\
%          Li et al.~\cite{li2019exploiting}&x&x& Build enc./dec. oracle \\
%          Li et al.~\cite{li2022systematic}& \\
%          \bottomrule
%     \end{tabularx}
%     \caption{Previously proposed attacks on SEV that requires the NPF controlled-channel}
%     \label{tab:npf-attack}
% \end{table}


The hypervisor can trigger \gls{npf} on any page mapped to a \gls{cvm} by unsetting the present bit (\texttt{p}-bit) in the \gls{npt}.
Later, when the victim \gls{cvm} accesses the targetted pages, a \texttt{\#NPF} is generated by the CPU, and the faulting address is written into the \gls{vmcb}.

Hence, by unsetting the \texttt{p}-bit of every page used by a \gls{cvm}, the hypervisor can monitor its memory accesses and code execution at the page level. By combining this coarse-grained information with prior knowledge about the target \gls{vm}, attackers can infer much more fine-grained information, such as the system call trace~\cite{hetzelt_security_2017}, or the sequence of functions that performs certain sensitive operations.


% For this reason, \thename focuses on mitigating the two most prominent side-channel attacks against \glspl{cvm}.


% There are two security impacts of the \gls{npf} controlled-channel.

% First, the hypervisor can peek into the \gls{cvm}'s execution states and potentially leak data.
% To this end, its impact is similar to the Page Fault controlled-channel on SGX~\cite{xu2015controlledchannel}, which has been shown to leak sensitive information.
% On SEV, researcher has been \gls{cvm}~\cite{werner2019severest}.
% TODO: We perform additional evaluations to confirm this.
% Second, leaking the execution states of the \gls{cvm} allows the attacker to employ more fine-grained attacks.
% We observe that most known attacks on \gls{sev} depend on the \gls{npf} side-channel.

% Li et al. \cite{li2019exploiting} uses the sequence of \gls{npf} to extract page accesses when serving SSH packet, then uses pattern matching to extract the page that contains the packet to be sent over SSH.

% \paragraph{Ciphertext Side-channels}
% The recent CipherLeak attack~\cite{li2021cipherleaks,li2022systematic} employs the \gls{npf} controlled-channel.


% \begin{table}[]
%     \centering
%     \begin{tabular}{ll}
%     \toprule
%          Side-channel & Leakage  \\
%          \midrule
%          Nested PF & Page-level read/write/execute access pattern  \\
%          Ciphertext & Cacheline-level write access pattern\\
%          \bottomrule
%     \end{tabular}
%     \caption{Caption}
%     \label{tab:my_label}
% \end{table}


\subsection{Unikernels}
Unikernels are specialized kernel images linked directly with the application to be executed on a single-address space environment. Due to them being highly specialized to the specific workload, applications running on unikernels achieve faster boot time, better I/O performance, and smaller image size~\cite{unikraft} compared to traditional kernels.


 \thename takes advantage of the absence of the user-kernel switch in unikernels to achieve efficient instrumentation on memory accesses.
We port an existing unikernel implementation, Unikraft~\cite{unikraft}, to work with AMD SEV-SNP to utilize its mature POSIX compatibility.


\subsection{Oblivious RAM (ORAM)}
\gls{oram} algorithms are well-studied and proven ways to eliminate discernable memory trace differences on a memory region.
Path ORAM~\cite{path-oram} is the most adapted algorithm due to its simplicity and efficiency.
Path \gls{oram} organizes data blocks into a tree structure. Accesses to data blocks within the trees are transformed into path access that seems random to an observer. 


% \thename utilizes an implementation of PathORAM~\cite{stefanov2018path} to manage the page pool.

\section{Design}
\label{wtg}
We are interested in two key features of the OS, (1) being able to control the virtual address mapping through page table manipulation, and (2) being able to intercept access to memory through page faults.
Through page table manipulation, we observe that we can control application \emph{access pattern} to physical memory without changing the \emph{virtual address space} of the application.

On the other hand, page fault-based interception enables a primitive similar to that of compiler instrumentation that inserts inline reference monitor before/after memory accesses~\cite{asan,wichelmann_cipherfix_nodate,dynpta} to enforce security policies. 
With access to the page table, we can clear the present bit in the page of interest such that all access to those pages would transfer to the kernel reference monitor.
However, two issues prevent the page fault handler from being a reliable defense tool.
First, page-level interception is too coarse-grained. After one instruction is intercepted, the control must be passed back to the faulting instruction. After this, there is no way to intercept other memory accesses to the same page. 
Second, traditional page fault handlers do not need to be aware of the semantics of the fault-inducting instruction, e.g., whether it is a load or store instruction and how many bytes are accessed. 
However, understanding instruction semantics is necessary to implement many fine-grained defense policies, like changing the \emph{location} and \emph{content} of data being accessed.
For the above reasons, we incorporate an \emph{emulation engine} that the faulting instruction and allows the fault handler to emulate the instruction. The emulation engine exposes hooks through callbacks that enable fine-grained security instrumentation on load and store instructions.

\textbf{Virtual Memory Areas}
Operating systems manage the address space through an abstraction called the \gls{vma}. A \gls{vma} contains the start and end address of a contiguous virtual address range. 
Especially in Unikraft, the infrastructure \gls{vma} is implemented as a library that allows the user to define \gls{vma} with user-defined behaviors.
We extend the \gls{vma} abstraction to fit our purposes and use it to implement our security mechanisms as \glspl{vma} within the address spaces. For instance, memory accesses virtual addresses that are within the oblivious paging \gls{vma} (\autoref{sec:ovma}) will have their paging pattern protected.





\subsection{Oblivious Paging}
\label{sec:ovma}
Our first defense tries to obfuscate application memory accesses to \emph{guest physical memory} to make them seem random and observing hypervisor. At a high level, we employ a single 2MB page as a trusted physical space (an ORAM stash) and enforce that accessible virtual addresses are only mapped to this page. We place the kernel and application memory into an ORAM tree, We use the Path ORAM algorithm to swap pages in and out of the stash. This way, an observing hypervisor that monitors page accesses would only see accesses to a single page, that is, the ORAM stash, along with seemingly random ORAM accesses.


\subsection{Automatic Salting of Sensitive Writes}
The recently discussed cipherleak attack~\cite{li_cipherleaks_nodate,li_systematic_2022} exploits the ciphertext side-channel to extract sensitive memory access patterns from cryptographic libraries. 
Such an attack can be prevented if the data is \emph{salted} with a nonce before being written to memory.
This approach is employed by a proposed defense that utilizes compiler analysis and binary instrumentation~\cite{wichelmann_cipherfix_nodate}.
With our emulation primitive, we can implement an automatic salting scheme without relying on a compiler pass.

% We implement an automatic salting scheme.

\subsection{In-register Data Protection}
The CacheWarp attack~\cite{cachewarp} selectively drops a sensitive memory write instruction from a \gls{cvm}, to enable control flow hijacking and bypassing of access control checks. 
Conceptually, this attack is similar to the \emph{data-only} attacks that corrupt sensitive in-memory data structures through memory vulnerabilities. 
To prevent data-only attacks, an approach is to protect sensitive data within a trusted memory region. Since all memory can be used for the attack, we use the vector registers. This approach can be seen in previous work that employs compiler instrumentation~\cite{dynpta}. 
We can implement such a defense without the compiler.


% test